\documentclass{handout}
\usepackage{handoutSetup}\usepackage{handoutShortcuts} 
\setboolean{BigAndWide}{true}
 
\begin{document}
\handoutHeader

\begin{verbatimwrite}{\jobname.title}
The Laibson Model of Time Inconsistency
\end{verbatimwrite}

\handoutNameMake

This handout provides a simple example of a discrete-time 
solution to the problem of a consumer with a self-control
problem {\it a la} \cite{laibson:goldeneggs}.\footnote{Laibson's
own lecture notes are available on his website, and are very
good; as of this writing, they are in lecture 6 of his course 2010c, 
handouts, lecture06.}

\newcommand{\bc}{\mathbf{c}}
\newcommand{\bu}{\uFunc}
\newcommand{\bv}{\vFunc}
\newcommand{\gv}{\mathfrak{v}}
\newcommand{\gc}{\mathfrak{c}}
\newcommand{\bw}{\mathbf{w}}


Suppose a value function $\bv_{t+1}(m_{t+1})$ exists for period $t+1$.  
Then for any period-$t$ consumption function $\bc_{t}$ we can define 
\begin{equation}\begin{gathered}\begin{aligned}
   {\bv}_{t}(m_{t};\pmb{\chi}_{t}) & =  \bu(\pmb{\chi}_{t}(m_{t})) + \phantom{\delta} \beta  \Ex_{t}[{\bv}_{t+1}((m_{t}-\pmb{\chi}_{t}(m_{t}))\Rfree+y_{t+1})] \label{eq:bv}
\\ {\gv}_{t}(m_{t};\pmb{\chi}_{t}) & =  \bu(\pmb{\chi}_{t}(m_{t})) +          \delta  \beta  \Ex_{t}[{\bv}_{t+1}((m_{t}-\pmb{\chi}_{t}(m_{t}))\Rfree+y_{t+1})] %\label{eq:gv}
\end{aligned}\end{gathered}\end{equation}

Notice that these functions are well defined for any consumption function $\pmb{\chi}_{t}(m_{t})$ 
that is feasible; they are {\it not} Bellman equations because they do not assume that
the consumption function $\pmb{\chi}_{t}$ is optimal.  For example, these functions would be
well defined for $\pmb{\chi}_{t}(m_{t}) = m_{t}$, or for $\pmb{\chi}_{t}(m_{t})=1$, or for many other
potential consumption rules.  

What these functions capture is the value of behaving according to the rule $\pmb{\chi}_{t}$ 
in the current period, under two possible assumptions about discounting of the future:
Either next period's value is discounted by the factor $\beta$ (for $\bv_{t}$) or by $\delta \beta$ (for
$\gv_{t}$).  

Now consider two possible candidates for $\pmb{\chi}_{t}$:
\begin{equation}\begin{gathered}\begin{aligned}
    \bc_{t}(m_{t}) & =  \argmax_{c}~~ \bu(c) +  \phantom{\delta} \beta \Ex_{t}[\bv_{t+1}((m_{t}-c)\Rfree+y_{t+1})] \label{eq:cexpvshyp}
\\  \gc_{t}(m_{t}) & =  \argmax_{c}~~ \bu(c) +           \delta  \beta \Ex_{t}[\bv_{t+1}((m_{t}-c)\Rfree+y_{t+1})] %\label{eq:chyp}
\end{aligned}\end{gathered}\end{equation}

If we solve the problem recursively using $\pmb{\chi}=\bc$ in every period, we
obtain the standard time consistent solution.  (Think about why).

The Laibson alternative is to suppose that there is something special
about ``now'': Next period's value is discounted not only by the standard
geometric discount factor $\beta$, but also by an extra factor $\delta$ (Laibson
argues that at an annual frequency the appropriate value of $\delta$ is about 0.7).
This may reflect the fact that certain areas of the brain associated with emotional
rewards are activated only by instant gratification, and are not activated by thoughts
of future gratification (see, e.g.,~\cite{laibsonNeuro}).

It is clear from comparing the equations in \eqref{eq:cexpvshyp} that the consumer 
with Laibson preferences will consume more in the current period, because he 
values future rewards less.  

More insight about the solution can be obtained from the modified Euler 
equation that can be derived for the Laibson problem.  This is derived as follows.

Note first that if $\pmb{\chi}_{t} = \gc_{t}$ the \handoutC{Envelope} theorem implies
that 
\begin{equation}\begin{gathered}\begin{aligned}
  \gv_{t}^{m}(m_{t}) & =  \bu^{\prime}(c_{t})
\end{aligned}\end{gathered}\end{equation}
while the first order condition from the maximization problem implies that
\begin{equation}\begin{gathered}\begin{aligned}
  \bu^{\prime}(c_{t}) & =  \delta \Rfree\beta \Ex_{t}[\bv_{t+1}^{m}(m_{t+1})] \label{eq:foc}
\\ & =  \delta \bv_{t}^{m}(m_{t};\gc_{t})
\end{aligned}\end{gathered}\end{equation}

Now note that for $\pmb{\chi}_{t}=\gc_{t}$ there is a simple identity linking $\bv$ and $\gv$:
\begin{equation}\begin{gathered}\begin{aligned}
  \delta \bv_{t} & =  \gv_{t} - (1-\delta) \bu(\gc_{t}(m_{t})) \label{eq:dbv}
\end{aligned}\end{gathered}\end{equation}
(to see this, multiply the first equation in \eqref{eq:bv} by $\delta$ and note that the difference between
the result and the second equation is $(1-\delta) \bu(c_{t})$).  Now differentiate \eqref{eq:dbv}
\begin{equation}\begin{gathered}\begin{aligned}
  \delta \bv_{t}^{m} & =  \gv_{t}^{m} - (1-\delta) \bu^{\prime}(c_{t})\gc_{t}^{m}(m_{t})
\end{aligned}\end{gathered}\end{equation}

Thus, 
\begin{equation}\begin{gathered}\begin{aligned}
  \uFunc^{\prime}(c_{t}) & =  \gv_{t}^{m}(m_{t};\pmb{\chi}_{t}) - (1-\delta) \bu^{\prime}(c_{t})\gc_{t}^{m}(m_{t}) \label{eq:LaibEul}
\end{aligned}\end{gathered}\end{equation}


If $\delta=1$, this collapses to the usual consumption Euler equation.  
However, if $\delta < 1$ (the Laibson case), the equation says 
several interesting things.  First, note that since $(1-\delta)$ 
and $\bu^{\prime}(c_{t})$ and $\gc^{m}_{t}$ are all positive, 
the contribution of the ``Laibson'' term in \eqref{eq:LaibEul} is
to reduce the RHS of the equation.  In order to match a lower RHS,
the LHS must be smaller.  But a smaller marginal utility of consumption
implies a higher level of consumption - so the Laibson consumer spends
more.

Second, notice that the magnitude of the ``present-bias'' effect depends
on the size of next period's marginal propensity to consume $\gc_{t}^{m}$.
If the MPC is small, the size of the Laibson bias will be small.

Finally, notice that this model nicely captures the commonplace psychological 
tension in which the cost of deviating from the optimal plan in a single period
may be trivially small (``eating dessert this one time will not make me fat''),
but the consequences of perpetual deviation could be quite large (``but if I give 
in to temptation this time, maybe that means I will always give in.'')  

 
\input handoutBibMake

\end{document}
